\documentclass[11pt]{amsart}
\usepackage{amsmath,amssymb}
\setlength{\emergencystretch}{3em}
\tolerance=600
\usepackage{tikz-cd}
\usepackage{tikz}
\usepackage[hidelinks]{hyperref}

\newtheorem{thm}{Theorem}[section]
\newtheorem{lem}[thm]{Lemma}
\newtheorem{cor}[thm]{Corollary}
\newtheorem{prop}[thm]{Proposition}
\theoremstyle{definition}
\newtheorem{defi}[thm]{Definition}
\newtheorem{rem}[thm]{Remark}
\newtheorem{ex}[thm]{Example}
\newtheorem{que}[thm]{Question}

\newcommand{\Set}{\mathbf{Set}}
\newcommand{\Fin}{\mathbf{FinSet}}
\newcommand{\Sub}{\mathrm{Sub}}
\newcommand{\Mod}{\mathrm{Mod}}
\newcommand{\Syn}{\mathrm{Syn}}
\newcommand{\PP}{\mathsf P}
\newcommand{\HH}{\mathsf H}
\newcommand{\cX}{\mathcal X}
\newcommand{\cC}{\mathcal C}
\newcommand{\cD}{\mathcal D}
\newcommand{\cE}{\mathcal E}
\newcommand{\cM}{\mathcal M}
\newcommand{\bT}{\mathbb T}
\newcommand{\id}{\mathrm{id}}
\newcommand{\cl}{\mathrm{cl}}

\begin{document}

\title[Intermediate logics and theories with finite sorts]{Intermediate
logics and the imaginaries of a theory whose models are finite on
every sort}
\author{Ryota Shibusawa}
\thanks{Corresponding author. E-mail: \texttt{r-shibusawa@daiichi-koudai.ac.jp}}
\address{Daiichi Institute of Technology}
\email{r-shibusawa@daiichi-koudai.ac.jp}
\keywords{intermediate logic, weak excluded middle, G\"odel--Dummett logic, Heyting category, pretopos completion, imaginaries, Kripke model, EI-category, coherent theory}
\subjclass[2020]{03G30 (primary), 03B55, 03C13, 18F10, 18B25 (secondary)}

\begin{abstract}
For an intermediate logic $L$ and a coherent theory, reading the
theory in first-order logic over $L$ and adding the imaginaries are two
operations which need not commute. For the coherent theories all of
whose models are finite on every sort we determine when they commute,
in terms of the category $\cM$ of models, except for the logics
proving linearity with a chain of three elements as a frame. For classical logic they
always commute. If $L$ does not prove the weak excluded middle, they
commute exactly when $\cM$ is a groupoid. If $L$ proves the weak
excluded middle and not the linearity axiom, they commute exactly when
every endomorphism of a model is invertible and no two non-invertible
homomorphisms are composable. For the logic of the Kripke frame $0<1$
they commute exactly when every endomorphism of a model is invertible.
This contains the known cases of finite categories of models and of
propositional theories. For the sufficiency the
theory is written as a directed colimit of theories in finite
languages, each with a finite category of models, and the
invertibility of the endomorphisms passes to cofinally many of them. The
necessity rests on an obstruction for all coherent theories: a
homomorphism which is injective on all sorts, the imaginary ones
included, and not surjective on a sort which is finite in its source
excludes commutation for every non-classical logic.
\end{abstract}

\maketitle

\section{Introduction}

For an intermediate logic $L$ and a small coherent category $\cC$ let
$\HH_L\cC$ be the syntactic category of the internal coherent theory of
$\cC$ in first-order logic over $L$, and $\PP$ the pretopos completion,
which adds the imaginaries. There is a full and faithful comparison
$c^L_\cC\colon\HH_L\PP\cC\to\PP\HH_L\cC$, and it is an equivalence if
and only if $\HH_L\PP\cC$ is a pretopos \cite[Lemma~2.2]{Interm}. We
say that \emph{the translation into $L$ commutes with the imaginaries
of $\cC$} if $c^L_\cC$ is an equivalence.

A \emph{model} of $\cC$ is a coherent functor $M\colon\cC\to\Set$, a
\emph{homomorphism} is a natural transformation, and $\Mod(\cC)$ is
the category of models. We say that $\cC$ \emph{has finite sorts} if
$M(A)$ is a finite set for every model $M$ and every object $A$. For
the syntactic category of a coherent theory this means that every
sort is finite in every model. Examples are the pretoposes $\Fin^{\cX}$
of the functors from a finite category $\cX$ to finite sets, whose
category of models is equivalent to $\cX$ if idempotents split in
$\cX$
\cite[\S7]{Heyting}; the bounded distributive lattices, that is, the
propositional theories, whose models are the prime filters; and the
theories of Example~\ref{ex:new}. A category is an \emph{EI-category}
if all its endomorphisms are invertible. The \emph{weak excluded
middle} is $\neg p\vee\neg\neg p$ and the \emph{linearity axiom} is
$(p\to q)\vee(q\to p)$; for $L$ proving linearity, $n_L\in\{1,2,\dots,
\infty\}$ is the supremum of the numbers of elements of the finite
chains which are frames of $L$, so that $n_L=1$ for classical logic
only and $n_L=2$ for the logic of the Kripke frame $0<1$.

\begin{thm}\label{thm:main}
Let $\cC$ be a small coherent category with finite sorts, $\cM=
\Mod(\cC)$, and $L$ an intermediate logic.
\begin{enumerate}
\item[(1)] If $L$ is classical logic, $c^L_\cC$ is an equivalence.
\item[(2)] If $L$ does not prove the weak excluded middle, $c^L_\cC$ is
an equivalence if and only if $\cM$ is a groupoid.
\item[(3)] If $L$ proves the weak excluded middle and not the
linearity axiom, $c^L_\cC$ is an equivalence if and only if $\cM$ is
an EI-category without two composable non-invertible homomorphisms.
\item[(4)] If $L$ proves linearity and $n_L=2$, $c^L_\cC$ is an
equivalence if and only if $\cM$ is an EI-category.
\item[(5)] If $L$ proves linearity and $n_L\ge3$: if $c^L_\cC$ is an
equivalence, $\cM$ is an EI-category; and if $\cM$ is an EI-category
without two composable non-invertible homomorphisms, $c^L_\cC$ is an
equivalence.
\end{enumerate}
\end{thm}

\begin{center}
\begin{tabular}{lll}
$L$&a frame of $L$&$c^L_\cC$ is an equivalence\\\hline
classical&the point&always\\
no w.e.m.&the fork&iff $\cM$ is a groupoid\\
w.e.m., not linear&the diamond&iff $\cM$ is EI, no two composable\\
&&non-invertible homomorphisms\\
linear, $n_L=2$&the chain $0<1$&iff $\cM$ is EI\\
linear, $n_L\ge3$&chains, $\le n_L$ elements&only if $\cM$ is EI; if
$\cM$ is EI, no\\
&&two composable non-invertible\\
&&homomorphisms
\end{tabular}
\end{center}
Here w.e.m.\ is the weak excluded middle. For $\cC=\Fin^{\cX}$ the
cases (1)--(4) are \cite[Thm.~1.1]{Classify}, where the condition in
case (5) is determined as well \cite{Strings}, and for a bounded
distributive lattice case (3) is \cite[Thm.~5.1]{Classify} and case
(4) is contained in \cite[Thm.~6.3]{Interm}.

The necessity in (3)--(5) uses an obstruction which does not need
finite sorts in all models.

\begin{thm}\label{thm:obstruction}
Let $\cC$ be a pretopos and $L$ an intermediate logic other than
classical logic. Suppose that there are a homomorphism $i\colon k\to m$
of models of $\cC$ which is injective on every object, and an object
$Y$ such that $Y(k)$ is finite and $Y(i)$ is not surjective. Then
$c^L_\cC$ is not an equivalence.
\end{thm}

A split monomorphism of models is injective on every object. So a
model with a finite sort which is a retract of a model in which this
sort is larger excludes the commutation for all non-classical logics:
for instance the sets with one and two elements for the theory of
pure equality, which gives \cite[Thm.~7.2]{Interm} with another proof,
or the groups with one and two elements for the theory of groups.

\subsection*{The method}
Under the hypothesis of finite sorts, a coherent type realised in one
model is realised in another one as soon as each of its formulas is
(Lemma~\ref{lem:types}); so a homomorphism between two models exists
if the sentences true in the first are true in the second, and in the
EI case the non-invertible homomorphisms are exactly those which
enlarge the set of true sentences (Proposition~\ref{prop:EI}). For
the sufficiency in (4) we consider, for every finite set $T$ of sorts
and relations of $\cC$, the theory $\bT_T$ of the coherent sequents
in this language which are valid in $\cC$. It has finitely many
models up to isomorphism, all of them finite, its pretopos is
equivalent to $\Fin^{\cX_T}$ for the finite category $\cX_T$ of its
models (Proposition~\ref{prop:finite}), and $\cC$ is the directed
colimit of these pretoposes. The models of $\bT_T$ need not be
restrictions of models of $\cC$, and $\cX_T$ need not be an
EI-category when $\cM$ is: for the pretopos of finite sets and a
language with one sort with two elements and no relations, the
constant maps are endomorphisms. But every finite language is
contained in one for which $\cX_T$ is an EI-category
(Proposition~\ref{prop:cofinal}), by an inverse limit argument. The result for finite categories of models
\cite{Strings} and the compatibility of $\HH_L$ with directed
colimits (Lemma~\ref{lem:colimit}) finish the proof.
\[
\begin{tikzcd}[column sep=large]
\Fin^{\cX_T}\simeq\cC_T\ar[r,"j_{TT'}"]\ar[rr,bend right=18,"j_T"']&
\cC_{T'}\simeq\Fin^{\cX_{T'}}\ar[r,"j_{T'}"]&\cC=\varinjlim_T\cC_T
\end{tikzcd}
\]

\subsection*{What is new and what is known}
Known: case (2), for every coherent category, $c^L$ being an
equivalence exactly in the Boolean case \cite[Thm.~1.1(1)]{Interm};
the obstruction over the diamond \cite[Prop.~5.1]{Interm}; the results
for $\Fin^{\cX}$ \cite{Strings,Classify} and for propositional
theories \cite{Interm,Classify} quoted above;
Theorem~\ref{thm:obstruction} for $\cC=\Fin^{\cX}$ and a split
monomorphism $i$ \cite[Thm.~3.1]{Classify}. New, to the author's
knowledge: Theorem~\ref{thm:obstruction} in general, and
Theorem~\ref{thm:main} beyond the two classes of examples mentioned.
Proposition~\ref{prop:finite}, that a coherent theory in a finite
language with bounded sorts is the theory of the finite functors on
its finite category of models, is presumably known; see
\cite[Cor.~2.6]{Ye24} for the pretoposes with an essentially finite
category of models. We include a direct proof. For classical logic the
comparison is governed, for all coherent theories, by a Galois
criterion \cite[Thm.~6.1]{Morley}; case (1) says that it is satisfied
by the theories with finite sorts, and our proof does not use the
criterion. A duality between
profinite monoids and the pretoposes of the coherent theories which
are categorical is given in \cite{Ye26}; we do not use it.

\section{The setting}\label{sec:setting}

Conventions are those of \cite{Interm}. An \emph{$L$-category} is a
Heyting category whose lattices of subobjects are Heyting algebras in
which the formulas of $L$ are valid; $\HH_L\cC$ is the free
$L$-category on the coherent category $\cC$, with unit $\eta$, which
we omit. Every object of $\HH_L\cC$ is a subobject $U\le A$ of an
object $A$ of $\cC$, and these subobjects are obtained from those in
$\cC$ by finite meets and joins, $\Rightarrow$, and $f^*$, $\exists_f$,
$\forall_f$ for the morphisms $f$ of $\cC$. The models of $\PP\cC$ are
the models of $\cC$, and $\PP\cC$ has finite sorts if $\cC$ has, an
object of $\PP\cC$ being a quotient of a finite coproduct of objects
of $\cC$; and $c^L_\cC$ is an equivalence if and only if $c^L_{\PP\cC}$
is. So \emph{from now on $\cC$ is a pretopos}; then $c^L_\cC$ is the
unit $\HH_L\cC\to\PP\HH_L\cC$, and it is an equivalence if and only if
$\HH_L\cC$ is a pretopos \cite[Lemma~2.2]{Interm}. A \emph{sentence}
is a subobject of the terminal object $1$ of $\cC$; it is \emph{true}
in a model $M$ if its value is not empty. The values of a model may be
empty sets, and the compactness and completeness theorems are used for
first-order logic with possibly empty sorts. Directed colimits of
categories are taken along functors which are compatible up to
coherent isomorphisms.

A \emph{model} of $\HH_L\cC$ is a coherent functor $\hat G\colon
\HH_L\cC\to\Set$; it is not required to be a Heyting functor. Its
restriction to $\cC$ is a model $M$ of $\cC$; we say that $\hat G$ is
\emph{over} $M$ and regard $\hat G(U)$, for $U\le A$, as a subset of
$M(A)$, which we also write $A(M)$. A natural transformation $\hat G\to
\hat G'$ of models over $M$, $M'$ restricts to a homomorphism
$\varphi\colon M\to M'$ and is determined by it; so a
\emph{homomorphism} $\hat G\to\hat G'$ is a homomorphism $\varphi\colon
M\to M'$ with $\varphi(\hat G(U))\subseteq\hat G'(U)$ for all $U$.

\begin{lem}[truth lemma]\label{lem:truth}
Let $\hat G$ be a model of $\HH_L\cC$ over $M$, $f\colon B\to A$ a
morphism of $\cC$, $U,V\le B$ in $\HH_L\cC$, $\alpha\in A(M)$ and
$\beta\in B(M)$.
\begin{enumerate}
\item[(a)] $\alpha\in\hat G(\forall_fU)$ if and only if for every
homomorphism $\varphi\colon\hat G\to\hat G'$, over $M\to M'$, and
every $\beta'\in B(M')$ with $f(\beta')=\varphi(\alpha)$ one has
$\beta'\in\hat G'(U)$.
\item[(b)] $\beta\in\hat G(U\Rightarrow V)$ if and only if for every
homomorphism $\varphi\colon\hat G\to\hat G'$, $\varphi(\beta)\in
\hat G'(U)$ implies $\varphi(\beta)\in\hat G'(V)$.
\end{enumerate}
\end{lem}
\begin{proof}
This is \cite[Lemma~2.2]{Strings}, stated there for $\cC=\Fin^{\cX}$;
we repeat the proof. The implications from left to right hold because
$\varphi$ preserves $\forall_fU$ and $U\Rightarrow V$, and
$f^*\forall_fU\le U$, $(U\Rightarrow V)\wedge U\le V$. For the
converse of (a) let $\alpha\notin\hat G(\forall_fU)$, and consider the
classical first-order theory of the coherent functors on $\HH_L\cC$
with a constant $c_b$ for every element $b$ of the sets $\hat G(X)$,
$X$ an object of $\HH_L\cC$, the axioms $g(c_b)=c_{g(b)}$ for the
morphisms $g$ of $\HH_L\cC$, a new constant $d$ of sort $B$ and the
axioms $f(d)=c_\alpha$ and $\neg U(d)$. If it had no model, then by
the compactness theorem and the completeness theorem for the coherent
category $\HH_L\cC$ there would be an object $X$ of $\HH_L\cC$, a
morphism $g\colon X\to A$ and $b\in\hat G(X)$ with $g(b)=\alpha$ such
that $f^*\exists_gX\le U$; then $V_0:=\exists_gX\le A$ satisfies
$\alpha\in\hat G(V_0)$ and $f^*V_0\le U$, so
$V_0\le\forall_fU$, a contradiction. A model of the theory is a
coherent functor $\hat G'$ with a natural transformation $\hat G\to
\hat G'$ and an element $\beta'\in B(M')\setminus\hat G'(U)$ over the
image of $\alpha$. (b) is the same, with the axioms $U(c_\beta)$ and
$\neg V(c_\beta)$.
\end{proof}

So the models of $\HH_L\cC$ jointly reflect the order of subobjects,
by the completeness theorem of coherent logic.

\subsection*{Kripke models}
A finite poset $K$ is a \emph{frame of $L$} if the Heyting algebra of
its up-sets is an $L$-algebra. Let $K$ be a finite poset such that
every $\uparrow\!k$ is a frame of $L$, and $M\colon K\to\Mod(\cC)$ a
functor. Then $\Set^K$ is an $L$-category and $M$ extends to a Heyting
functor $\HH_L\cC\to\Set^K$, whose value at $U\le A$ is the family of
the subsets of $A(M_k)$ defined by the forcing relation of the Kripke
model $M$ \cite[Lemma~2.3]{Interm}. The evaluation at a node $k$ is a
model of $\HH_L\cC$ over $M_k$, and the transitions are homomorphisms
between these models. For the one-point poset this is the
\emph{classical model} $\cl_M$ over $M$, the Heyting functor
$\HH_L\cC\to\Set$ extending $M$; for an isomorphism $\psi\colon M\to
M'$ we have $\psi(\cl_M(U))=\cl_{M'}(U)$. The two-element chain
$0<1$ is a frame of every intermediate logic other than classical
logic \cite[Lemma~7.1]{Interm}; so for such $L$ and a homomorphism
$a\colon M\to N$ we have the Kripke model with $M_0=M$, $M_1=N$ and
the transition $a$, and we write $\hat G_a$ for its model at the node
$0$. It is a model over $M$, the node $1$ carries $\cl_N$, and
$a\colon\hat G_a\to\cl_N$ is a homomorphism. The fork is a frame of
$L$ if and only if $L$ does not prove the weak excluded middle; if $L$
proves it and not linearity, the diamond is a frame of $L$
\cite[Lemma~3.1]{Interm}.
\[
\begin{tikzcd}[column sep=small,row sep=small]
\bullet&&\bullet&&&&\bullet&\\
&\bullet\ar[ul]\ar[ur]&&&&\bullet\ar[ur]&&\bullet\ar[ul]\\
&&&&&&\bullet\ar[ul]\ar[ur]&\\[-1ex]
&\text{fork}&&&&&\text{diamond}&
\end{tikzcd}
\]

\section{The obstruction}

\begin{proof}[Proof of Theorem~\ref{thm:obstruction}]
Let $c_1,\dots,c_N$ be the elements of $Y(k)$, $u=(c_1,\dots,c_N)\in
Y^N(k)$, and
\[
Q:=\exists_p\bigwedge_{j=1}^N\neg\,[z=x_j]\ \le\ Y^N
\]
in $\HH_L\cC$, where $[z=x_j]\le Y\times Y^N$ is the subobject of the
$(z,x_1,\dots,x_N)$ with $z=x_j$ and $p\colon Y\times Y^N\to Y^N$ is
the projection. We use the homomorphism $i\colon\hat G_i\to\cl_m$ of
models of $\HH_L\cC$. First, $i(u)\in\cl_m(Q)$: an element $z$ of
$Y(m)$ outside the image of $Y(i)$ is different from all $i(c_j)$, and
$\cl_m$ is a Heyting functor into $\Set$. Second, $u\notin\hat G_i(Q)$:
every $z\in Y(k)$ is some $c_j$, and then the node $0$ does not force
$\neg[z=x_j]$ at $(z,u)$.

Let $X=Y^N\sqcup Y^N$ with the inclusions $j_1,j_2$, and $E\le X\times
X$ the join of the diagonal and of the images of $Q$ under $(j_1,j_2)$
and $(j_2,j_1)$, an equivalence relation in $\HH_L\cC$. If $\HH_L\cC$
were a pretopos, the quotient $X/E$ would be a subobject of some
object $A$ of $\cC$, and the map $\hat G_i(X/E)\to\cl_m(X/E)$ induced
by the homomorphism $i$ would be a restriction of $A(i)$, which is
injective. But models preserve quotients of equivalence relations; the
classes of $j_1(u)$ and $j_2(u)$ are different in $\hat G_i(X/E)$ and
are sent to the classes of $j_1(i(u))$ and $j_2(i(u))$, which are
equal in $\cl_m(X/E)$.
\end{proof}

The injectivity is required on all objects of the pretopos. A
homomorphism which is injective on the sorts of a theory need not be
injective on the imaginary sorts: for the theory of a sort with a
relation $\ne$ which is the complement of the equality, the
homomorphisms are the injective maps, and the inclusion of a set with
two elements in a set with three elements is not injective on the
quotient of the sort by the equivalence relation ``$x=y$, or there
are three different elements''.

\section{Theories with finite sorts}\label{sec:finite}

Let $\cC$ be a coherent category.

\begin{lem}\label{lem:bounded}
If $A(M)$ is finite for every model $M$, there is $n$ such that $A(M)$
has at most $n$ elements for every $M$; equivalently, the join of the
diagonals $[x_i=x_j]$, $i<j$, is the top subobject of $A^{n+1}$.
\end{lem}
\begin{proof}
Otherwise the classical theory of the models of $\cC$ with infinitely
many new constants of sort $A$ and the axioms that they are different
is finitely satisfiable, and a model of it is a model of $\cC$ in
which $A$ is infinite. The equivalence is the completeness theorem.
\end{proof}

\begin{lem}\label{lem:retract}
Let $\varepsilon$ be an idempotent endomorphism of a model $M$. Then
the sets $k(A)=\{x\in A(M)\mid\varepsilon(x)=x\}$ form a model $k$,
the inclusion $i\colon k\to M$ is a homomorphism, and
$\varepsilon=ir$ for a homomorphism $r\colon M\to k$ with $ri=\id$.
\end{lem}
\begin{proof}
$k$ is a subfunctor of $M$ by naturality, and a retract of $M$ in the
category of functors by $r=\varepsilon$. A retract of an isomorphism
is an isomorphism, so $k$ preserves finite limits; a retract of a
surjection is a surjection, so $k$ preserves covers; and for
subobjects $S,S'\le A$ an element $x\in k(S\vee S')$ lies in $S(M)$,
say, and then $x=\varepsilon(x)\in k(S)$, while $k(0)\subseteq
M(0)=\emptyset$. So $k$ preserves finite joins of subobjects.
\end{proof}

We use repeatedly that an inverse system of non-empty finite sets over
a directed set has a non-empty limit. The next lemma replaces, for
finite sorts, the realisation of types
in saturated models.

\begin{lem}[types]\label{lem:types}
Let $N$, $N'$ be models of $\cC$ which are finite on every object, and
$(\alpha_i)_{i\in I}$, $(\alpha'_i)_{i\in I}$ families of elements
$\alpha_i\in A_i(N)$, $\alpha'_i\in A_i(N')$. Suppose that for all
$i_1,\dots,i_r\in I$, $r\ge0$, and every subobject $S\le A_{i_1}\times\dots\times
A_{i_r}$ in $\cC$,
\[
(\alpha_{i_1},\dots,\alpha_{i_r})\in S(N)\quad\text{implies}\quad
(\alpha'_{i_1},\dots,\alpha'_{i_r})\in S(N').
\]
Then there is a homomorphism $\lambda\colon N\to N'$ with
$\lambda(\alpha_i)=\alpha'_i$ for all $i$. (For $r=0$ the hypothesis
says that the sentences true in $N$ are true in $N'$.)
\end{lem}
\begin{proof}
For a finite set $D$ of objects, a finite set $F$ of morphisms between
objects of $D$ and a finite set $I_0\subseteq I$ with $A_i\in D$ for
$i\in I_0$, let $\Lambda(D,F,I_0)$ be the set of the families of maps
$\lambda_A\colon A(N)\to A(N')$, $A\in D$, which commute with the
morphisms in $F$ and satisfy $\lambda(\alpha_i)=\alpha'_i$ for $i\in
I_0$. These are finite sets, they form an inverse system over a
directed set, and an element of the inverse limit is a natural
transformation $N\to N'$, a homomorphism, with the required property.
So it suffices to show that $\Lambda(D,F,I_0)$ is not empty.

Let $\nu$ be the tuple of all elements of the sets $A(N)$, $A\in D$,
without repetitions, an element of $B(N)$ for a finite product $B$ of
objects of $D$, and $\alpha=(\alpha_i)_{i\in I_0}\in B_0(N)$, $B_0=
\prod_{i\in I_0}A_i$; likewise $\alpha'\in B_0(N')$. Let $p$ be the set
of the subobjects $\Sigma\le B_0\times B$ with $(\alpha,\nu)\in\Sigma
(N)$; it is closed under finite meets. For $\Sigma\in p$ the image
$\exists\Sigma\le B_0$ under the projection contains $\alpha$ in $N$,
so it contains $\alpha'$ in $N'$ by the hypothesis: there is $\nu'\in
B(N')$ with $(\alpha',\nu')\in\Sigma(N')$. As $B(N')$ is finite and $p$
is closed under finite meets, one $\nu'$ satisfies this for all
$\Sigma\in p$. Let $\lambda_A$ send the entries of $\nu$ in $A(N)$ to
the corresponding entries of $\nu'$. For $f\colon A\to A'$ in $F$ and
an entry $\nu_j$ in $A(N)$, the element $f(\nu_j)$ is an entry
$\nu_{j'}$, and the subobject of the $(x,y)$ with $f(y_j)=y_{j'}$
belongs to $p$; so $f(\nu'_j)=\nu'_{j'}$, and $\lambda$ commutes with
$f$. For $i\in I_0$ the element $\alpha_i$ is an entry $\nu_j$, the
subobject $[x_i=y_j]$ belongs to $p$, and so $\lambda(\alpha_i)=\nu'_j
=\alpha'_i$.
\end{proof}

For the empty family: if every sentence true in $N$ is true in $N'$,
there is a homomorphism $N\to N'$.

\begin{prop}\label{prop:EI}
Let $\cC$ be a pretopos with finite sorts.
\begin{enumerate}
\item[(a)] If an endomorphism $h$ of a model $M$ is not invertible,
there is a split monomorphism $i\colon k\to M$ of models which is not
surjective on some object.
\item[(b)] $\Mod(\cC)$ is an EI-category if and only if every
homomorphism which is injective on every object is invertible.
\item[(c)] If $\Mod(\cC)$ is an EI-category, a homomorphism $h\colon
M\to N$ is invertible if and only if every sentence true in $N$ is
true in $M$.
\end{enumerate}
\end{prop}
\begin{proof}
(a) For an object $A$ let $\varepsilon_A$ be the idempotent power of
the map $h_A$ of the finite set $A(M)$; the powers $h_A^e$ equal to
$\varepsilon_A$ are those with $e$ large and divisible by the period
of $h_A$. For a morphism $f\colon A\to B$ there is $e$ with $h_A^e=
\varepsilon_A$ and $h_B^e=\varepsilon_B$, so $\varepsilon$ is natural,
an idempotent endomorphism of $M$. If $\varepsilon$ is the identity,
every $h_A$ is bijective and $h$ is invertible. Otherwise
Lemma~\ref{lem:retract} gives $i\colon k\to M$, not surjective on an
object on which $\varepsilon$ is not the identity.

(c) Let every sentence true in $N$ be true in $M$. By
Lemma~\ref{lem:types} there is a homomorphism $g\colon N\to M$. The
endomorphisms $gh$ and $hg$ are invertible, so $h$ is. The converse is
clear.

(b) If the homomorphisms which are injective on every object are
invertible, $\Mod(\cC)$ is an EI-category by (a). Conversely let
$\Mod(\cC)$ be an EI-category and $i\colon k\to M$ injective on every
object. For a sentence $s\le1$ let $1\sqcup_s1$ be the quotient of
$1\sqcup1$ by the equivalence relation which identifies the two
points over $s$; it has one element in a model in which $s$ is true
and two elements otherwise. As $i$ is injective on $1\sqcup_s1$, a
sentence true in $M$ is true in $k$, and $i$ is invertible by (c).
\end{proof}

\begin{lem}\label{lem:boolean}
A pretopos $\cC$ with finite sorts is Boolean if and only if
$\Mod(\cC)$ is a groupoid.
\end{lem}
\begin{proof}
Let $\cC$ be Boolean and $h\colon M\to N$ a homomorphism. It preserves
complements, so it reflects the subobjects of $\cC$; in particular it
is injective on every object $A$. For $e\ge1$ let $E_e\le1$ be the
image of the complement in $A^e$ of the join of the diagonals; it is
true in a model iff $A$ has at least $e$ elements there. As $h$
preserves and reflects $E_e$, the sets $A(M)$ and $A(N)$ have the same
number of elements, and $h$ is bijective on $A$. Conversely, if a
subobject $S\le A$ has no complement, there are a homomorphism
$h\colon M\to N$ and $a\in A(M)\setminus S(M)$ with $h(a)\in S(N)$
\cite[Lemma~4.1]{Heyting}, and $h$ is not invertible.
\end{proof}

\section{Finite languages}\label{sec:approx}

\begin{prop}\label{prop:finite}
Let $\bT$ be a coherent theory in a language with finitely many sorts
and relation symbols and no function symbols, which proves for every
sort $A$ a sequent $\top\vdash\bigvee_{i<j\le n}x_i=x_j$ in variables
$x_0,\dots,x_n$ of sort $A$. Let $\cX$ be a skeleton of the category of
models of $\bT$. Then $\cX$ is a finite category in which idempotents
split, and the evaluation $\mathrm{ev}\colon\PP\Syn(\bT)\to\Fin^{\cX}$
is an equivalence.
\end{prop}
\begin{proof}
A model has at most $n$ elements in each sort, so there are finitely
many models up to isomorphism and finitely many homomorphisms between
two of them; so $\cX$ is finite, and
idempotents split by Lemma~\ref{lem:retract}. The models of the
pretopos $\cD=\PP\Syn(\bT)$ are those of $\bT$, the evaluation is a
coherent functor between pretoposes, and it reflects the order of
subobjects by the completeness theorem; in particular it is faithful,
and the lattices $\Sub_\cD(X)$ are finite, being embedded in the
lattices of subfunctors of $\mathrm{ev}(X)$.

Let $M$ be a model, $\nu$ the tuple of all its elements, without
repetitions, an element of $B(M)$ for a product $B$ of sorts, and
$S_M\le B$ the meet of the subobjects which contain $\nu$ in $M$. For
a model $N$ and $\nu'\in S_M(N)$ the map $\nu_j\mapsto\nu'_j$ preserves
the relation symbols, the atomic formulas true of $\nu$ being above
$S_M$; conversely a homomorphism $\lambda\colon M\to N$ gives
$\lambda(\nu)\in S_M(N)$. So $\mathrm{ev}(S_M)\cong\cX(M,-)$, with
$\nu$ corresponding to the identity.

\emph{The evaluation is full on subobjects.} Let $X$ be an object of
$\cD$, a quotient of a finite coproduct of subobjects $X_l$ of
products of sorts, $W\le\mathrm{ev}(X)$ a subfunctor and $w\in W(M)$.
Then $w$ is the class of a tuple in some $X_l(M)$, whose entries are
entries of $\nu$; so it is the image of $\nu$ under a morphism $t_w
\colon S_M\to X_l\to X$ of $\cD$ given by a projection, $S_M$ being
contained in the pullback of $X_l$ along this projection because
$\nu$ is. The image of $\mathrm{ev}(t_w)$ is the subfunctor of
$\mathrm{ev}(X)$ generated by $w$. Hence $W$ is the join of the images
of the $\mathrm{ev}(t_w)$ for finitely many $w$, and the evaluation
preserves images and joins.

\emph{The evaluation is full and essentially surjective.} The graph of
a morphism $\mathrm{ev}(X)\to\mathrm{ev}(X')$ is $\mathrm{ev}(\Gamma)$
for a subobject $\Gamma\le X\times X'$, which is the graph of a
morphism because the evaluation reflects the order of subobjects.
Every functor $W\colon\cX\to\Fin$ is the quotient of a finite
coproduct of representable functors, the value of $\mathrm{ev}$ at a
coproduct $Z$ of objects $S_M$, by an equivalence relation, which is
$\mathrm{ev}(E)$ for some $E\le Z\times Z$; $E$ is an equivalence
relation, and $\mathrm{ev}(Z/E)\cong W$.
\end{proof}

From now on $\cC$ is a pretopos with finite sorts. A \emph{finite
language} is a pair $T=(S,R)$ of a finite set $S$ of objects of $\cC$
and a finite set $R$ of subobjects of finite products of objects in
$S$; we write $T\subseteq T'$ if $S\subseteq S'$ and $R\subseteq R'$.
Let $\bT_T$ be the coherent theory with the sorts $S$, the relation
symbols $R$, and as axioms all coherent sequents in this language
which are valid in $\cC$. By Lemma~\ref{lem:bounded} it satisfies the
hypothesis of Proposition~\ref{prop:finite}. Let $\cX_T$ be a
skeleton of its category of models, $\cC_T=\PP\Syn(\bT_T)\simeq
\Fin^{\cX_T}$, and $j_T\colon\cC_T\to\cC$ the coherent functor
extending the interpretation of the formulas in $\cC$. For $T\subseteq
T'$ every model of $\bT_{T'}$ restricts to a model of $\bT_T$, and
there is a coherent functor $j_{TT'}\colon\cC_T\to\cC_{T'}$ with
$j_{T'}j_{TT'}\cong j_T$.

\begin{lem}\label{lem:union}
Each $j_T$ reflects the order of subobjects and is faithful. Every
object of $\cC$ is isomorphic to $j_T(X)$ for some $T$ and $X$, and
every morphism $j_T(X)\to j_T(X')$ of $\cC$ is of the form
$j_{T'}(f)$, up to the isomorphisms $j_{T'}j_{TT'}\cong j_T$, for some
$T'\supseteq T$ and some $f\colon j_{TT'}(X)\to j_{TT'}(X')$. So $\cC$
is the colimit of the directed system of the categories $\cC_T$, and
of every cofinal subsystem.
\end{lem}
\begin{proof}
The functor $\Syn(\bT_T)\to\cC$ reflects the order of subobjects by
the definition of $\bT_T$. The subobjects of a quotient $(\coprod_lX_l)
/F$ in a pretopos completion are the families of subobjects of the
$X_l$ which are saturated for $F$; so $j_T$ reflects the order of
subobjects, and a coherent functor with this property is faithful. An
object $A$ of $\cC$ is a sort of $T=(\{A\},\emptyset)$. Let $X=
(\coprod_lX_l)/F$ and $X'=(\coprod_{l'}X'_{l'})/F'$ in $\cC_T$ and
$g\colon j_T(X)\to j_T(X')$. The pullbacks $\Gamma_{ll'}\le j_T(X_l)
\times j_T(X'_{l'})$ of the graph of $g$ are subobjects of products of
objects in $S$; let $T'=(S,R\cup\{\Gamma_{ll'}\})$. In $\cC_{T'}$ the
relation symbols $\Gamma_{ll'}$ define a subobject $\Gamma$ of the
coproduct of the $X_l\times X'_{l'}$ which is saturated and is the
graph of a morphism $f$, because $j_{T'}$ reflects the order of
subobjects; and $j_{T'}(f)=g$.
\end{proof}

\begin{lem}[directed colimits]\label{lem:colimit}
Let a pretopos $\cC$ be the colimit of a directed system of pretoposes
$\cC_i$ and coherent functors. If $\HH_L\cC_i$ is a pretopos for all
$i$, then $\HH_L\cC$ is a pretopos.
\end{lem}
\begin{proof}
Let $\cD$ be the colimit of the categories $\HH_L\cC_i$ along the
Heyting functors $\HH_L(j_{ii'})$. A finite diagram in $\cD$ comes from
a diagram in some $\HH_L\cC_i$, and two parallel morphisms of
$\HH_L\cC_i$ which become equal in $\cD$ are equal in some
$\HH_L\cC_{i'}$. Hence $\cD$ has finite limits, computed at a stage; a
morphism of $\cD$ is a monomorphism iff it is represented by a
monomorphism at some stage, its kernel pair being computed there; so
$\Sub_\cD(X)$ is the directed colimit of the lattices of subobjects of
the representatives of $X$. Images, finite joins and the operations
$\forall_f$ are computed at a stage and are preserved by the
transitions; the adjunctions which characterise them, being
inequalities between finitely many subobjects, hold in $\cD$; and
$\Sub_\cD(X)$, a directed colimit of $L$-algebras and homomorphisms,
is an $L$-algebra. So $\cD$ is an $L$-category and the functors
$\HH_L\cC_i\to\cD$ are Heyting. In the same way, a functor from $\cD$
to an $L$-category is Heyting iff its restrictions to the $\HH_L\cC_i$
are, and a functor from $\cC$ to a coherent category is coherent iff
its restrictions to the $\cC_i$ are. Therefore, for an $L$-category
$\cE$, the Heyting functors $\cD\to\cE$ correspond to the compatible
families of Heyting functors $\HH_L\cC_i\to\cE$, these to the
compatible families of coherent functors $\cC_i\to\cE$, and these to
the coherent functors $\cC\to\cE$: $\cD$ is $\HH_L\cC$.

The transitions $\HH_L(j_{ii'})$ are coherent functors between
pretoposes, so they preserve finite coproducts and quotients of
equivalence relations. An equivalence relation in $\cD$ is represented
by an equivalence relation at some stage, the conditions being
inequalities between finitely many subobjects; its quotient at this
stage is a quotient in $\cD$, because a coequaliser which is preserved
by all transitions is a coequaliser in the colimit, and it is
effective and stable under pullback because this holds at every
stage. The same holds for finite coproducts. So $\cD$ is a pretopos.
\end{proof}

Since the pretoposes $\HH_L\cC_i$ for $i$ in a cofinal subset form a
system with the same colimit, it suffices in
Lemma~\ref{lem:colimit} that $\HH_L\cC_i$ is a pretopos for all $i$
in a cofinal subset.

\begin{prop}\label{prop:cofinal}
Let $\cC$ be a pretopos with finite sorts such that $\Mod(\cC)$ is an
EI-category. Then every finite language is contained in a finite
language $T$ such that $\cX_T$ is an EI-category.
\end{prop}
\begin{proof}
Let $S$ be a finite set of objects, and $\cC_S$ the full subcategory
of $\cC$ of the subobjects of the finite products of objects in $S$;
it is closed under finite limits, images and subobjects, so it is a
coherent category and the inclusion is coherent; for a model $M$ of
$\cC$ we write $M|_S$ for its restriction to $\cC_S$, which has more
objects than those in $S$; the models of $\cC_S$ are
finite on every object, the bounds of Lemma~\ref{lem:bounded} being
valid in $\cC_S$.

\emph{Step 1: every model $m$ of $\cC_S$ is isomorphic to the
restriction $M|_S$ of a model $M$ of $\cC$.} Consider the classical
theory of the models of $\cC$ with a constant for every element of
the sets $m(A)$, $A\in S$, and the axioms $\Sigma(c)$ or
$\neg\Sigma(c)$, for the objects $\Sigma$ of $\cC_S$ and the tuples
$c$ of constants, according to whether $c\in m(\Sigma)$ or not. A
finite part of it says $\Sigma(c)$ and $\neg\Sigma_1(c),\dots,
\neg\Sigma_e(c)$ for one tuple $c$ and objects of $\cC_S$ with
$c\in m(\Sigma)$ and $c\notin m(\Sigma_1\vee\dots\vee\Sigma_e)$; so
$\Sigma\not\le\Sigma_1\vee\dots\vee\Sigma_e$ in $\cC$, and there is a
model of $\cC$ with an element of $\Sigma$ outside the $\Sigma_d$. By
compactness there are a model $M$ of $\cC$ and a homomorphism $u\colon
m\to M|_S$ of models of $\cC_S$ which reflects the subobjects of
$\cC_S$. By Lemma~\ref{lem:types}, applied to $\cC_S$, the models
$M|_S$ and $m$, and the families of the elements $u(x)$ and $x$, for
$x$ in the sets $m(A)$, there is $r\colon M|_S\to m$ with $ru=\id$.
\[
\begin{tikzcd}
m\ar[r,"u"]\ar[dr,"\id"']&M|_S\ar[d,"r"]\ar[r,"ur"]&M|_S\\
&m\ar[ur,"u"']&
\end{tikzcd}
\qquad\qquad
\begin{tikzcd}
M\ar[r,dashed,"\lambda"]&M\\
M|_S\ar[r,"ur"]\ar[u,dash,dotted]&M|_S\ar[u,dash,dotted]
\end{tikzcd}
\]

Every homomorphism $g\colon M|_S\to M'|_S$, for models $M$, $M'$ of
$\cC$, is the restriction of a homomorphism $M\to M'$: apply
Lemma~\ref{lem:types} to $\cC$ and the families of the elements $x$
and $g(x)$, for $x$ in the sets $A(M)$, $A\in S$; a subobject in $\cC$
of a product of objects in $S$ is an object of $\cC_S$ and is
preserved by $g$. So the idempotent $ur$ of $M|_S$ is the restriction
of an endomorphism of $M$, which is invertible; hence $ur$ is
injective on every object of $\cC_S$, and an injective idempotent is
the identity. So $u$ is an isomorphism.

\emph{Step 2: $\Mod(\cC_S)$ is an EI-category.} By Step 1 an
endomorphism of a model of $\cC_S$ is, up to isomorphism, the
restriction of an endomorphism of a model of $\cC$.

\emph{Step 3.} Let $n_A$ be a bound for the object $A\in S$ as in
Lemma~\ref{lem:bounded}. For a finite set $R$ of objects of $\cC_S$
let $\mathcal E_R$ be the set of the pairs $(m,g)$ of a model $m$ of
$\bT_{(S,R)}$ whose sorts are sets $\{1,\dots,e_A\}$, $e_A\le n_A$,
and an endomorphism $g$ of $m$ which is not injective on some sort. It
is a finite set, and for $R\subseteq R'$ the restriction of the
structure is a map $\mathcal E_{R'}\to\mathcal E_R$. If all
$\mathcal E_R$ were non-empty, an element of the inverse limit would
be a structure $m$ interpreting all objects of $\cC_S$ as relations
on the sets $\{1,\dots,e_A\}$, in which all coherent sequents valid
in $\cC$ hold, with an endomorphism $g$ which is not injective. Since
the finite limits, the images and the finite joins of $\cC_S$, and the
graphs of its morphisms, are described by such sequents, $m$ is a
model of $\cC_S$ and $g$ a non-invertible endomorphism, against
Step 2. So $\mathcal E_{R_0}$ is empty for some $R_0$, and then
$\mathcal E_R$ is empty for all $R\supseteq R_0$.

\emph{Step 4.} Let $(S,R)$ be a finite language and $T=(S,R\cup R_0)$.
Every endomorphism $g$ of a model of $\bT_T$ is injective, hence
bijective, on every sort; its inverse is a power of $g$ and so a
homomorphism. So $\cX_T$ is an EI-category.
\end{proof}

\begin{thm}\label{thm:two}
Let $\cC$ be a pretopos with finite sorts.
\begin{enumerate}
\item[(a)] For classical logic, $c^L_\cC$ is an equivalence.
\item[(b)] Let $L$ prove linearity with $n_L=2$. Then $c^L_\cC$ is an
equivalence if and only if $\Mod(\cC)$ is an EI-category.
\end{enumerate}
\end{thm}
\begin{proof}
(a) For every finite language $T$ the category $\HH_L\cC_T\simeq\HH_L
\Fin^{\cX_T}$ is a pretopos \cite[Prop.~5.3]{Strings}, $\cX_T$ being a
finite category in which idempotents split. By Lemmas~\ref{lem:union}
and~\ref{lem:colimit}, $\HH_L\cC$ is a pretopos.

(b) If $\Mod(\cC)$ is an EI-category, the finite languages $T$ with
$\cX_T$ an EI-category are cofinal by Proposition~\ref{prop:cofinal};
for these $\HH_L\cC_T$ is a pretopos \cite[Cor.~6.3]{Strings}, and
Lemmas~\ref{lem:union} and~\ref{lem:colimit} apply. If $\Mod(\cC)$ is
not an EI-category, Proposition~\ref{prop:EI}(a) gives a homomorphism
$i\colon k\to M$ and an object $Y$ as in
Theorem~\ref{thm:obstruction}, the set $Y(k)$ being finite.
\end{proof}

The second half of this proof shows, for a pretopos with finite sorts
and every non-classical $L$: if $c^L_\cC$ is an equivalence,
$\Mod(\cC)$ is an EI-category.

\section{The weak excluded middle}

\begin{lem}\label{lem:classical}
Let $\cC$ be a pretopos, $L$ an intermediate logic and $\hat G$ a
model of $\HH_L\cC$ over $M$ such that every homomorphism out of
$\hat G$ has an invertible underlying homomorphism. Then $\hat G=
\cl_M$.
\end{lem}
\begin{proof}
If $\psi\colon\hat G\to\hat G'$ is a homomorphism, then $\hat G'$ has
the same property, a homomorphism $\lambda$ out of $\hat G'$ giving
the homomorphism $\lambda\psi$ out of $\hat G$. We show $\hat G(U)=
\cl_M(U)$ for all models with this property at once, by induction on
the construction of $U$. This is clear for the subobjects in $\cC$
and compatible with the coherent operations. By
Lemma~\ref{lem:truth}, $\alpha\in\hat G(\forall_fU)$ if and only if
$\beta'\in\hat G'(U)=\cl_{M'}(U)$ for all homomorphisms $\psi\colon
\hat G\to\hat G'$ and all $\beta'$ over $\psi(\alpha)$; since $\psi$
is invertible and the classical models are compatible with the
isomorphisms, this holds if and only if $\beta\in\cl_M(U)$ for all
$\beta$ over $\alpha$, the identity being among the $\psi$; that is,
if and only if $\alpha\in\cl_M(\forall_fU)$. The case of $\Rightarrow$
is the same.
\end{proof}

\begin{lem}\label{lem:top}
Let $L$ be an intermediate logic other than classical logic which
proves the weak excluded middle, and $\cC$ a pretopos with finite
sorts such that $\Mod(\cC)$ is an EI-category without two composable
non-invertible homomorphisms. Let $\hat G$ be a model of $\HH_L\cC$
over $M$, and $R$ the collection of the non-invertible homomorphisms
$a\colon M\to N$ of models of $\cC$ which are homomorphisms from
$\hat G$ to some model over $N$. Then either $R$ is empty and $\hat G=
\cl_M$, or $R$ consists of the homomorphisms $\iota a$, for one
$a\colon M\to N$ and all isomorphisms $\iota$ out of $N$, and
$\hat G=\hat G_a$.
\end{lem}
\begin{proof}
If $a\colon M\to N$ is not invertible, every homomorphism out of $N$
is invertible; so by Lemma~\ref{lem:classical} the target of a
homomorphism $a\in R$ is $\cl_N$. If $R$ is empty, $\hat G=\cl_M$ by
Lemma~\ref{lem:classical}.

Let $a\colon M\to N$ and $a'\colon M\to N'$ be in $R$. If $a'=\lambda
a$ for a homomorphism $\lambda$, then $\lambda$ is invertible. Suppose
that there is no such $\lambda$. By Lemma~\ref{lem:types}, applied to
the families of the elements $a(x)$ and $a'(x)$, for $x$ in the sets
$A(M)$, there are an object $A$, an element $\alpha\in A(M)$ and a
subobject $S\le A$ in $\cC$ with $a(\alpha)\in S(N)$ and
$a'(\alpha)\notin S(N')$; so $a'(\alpha)\in\cl_{N'}(\neg S)$. In
$\HH_L\cC$ the subobject $\neg S\vee\neg\neg S$ of $A$ is the top
element; so $\alpha\in\hat G(\neg S)$, and then $a(\alpha)\in
\cl_N(\neg S)$ by the homomorphism $a$, or $\alpha\in\hat G(\neg\neg
S)$, and then $a'(\alpha)\in\cl_{N'}(\neg\neg S)$ by $a'$. Both are
impossible. Conversely $\iota a\in R$ for $a\in R$ and an isomorphism
$\iota$, as $\iota\colon\cl_N\to\cl_{N'}$ is a homomorphism. So $R$,
if not empty, consists of the $\iota a$.

Let $R$ be non-empty, $a\in R$, and let $\psi\colon\hat G\to\hat G'$
be a homomorphism with $\psi\colon M\to M'$ invertible; let $R'$ be
the collection defined for $\hat G'$ as $R$ for $\hat G$. Then
$R'\psi\subseteq R$. We show that $R'$ is not empty. Otherwise
$\hat G'=\cl_{M'}$ by Lemma~\ref{lem:classical}. By
Proposition~\ref{prop:EI}(c) there is a sentence $s$ which is true in
$N$ and not in $M$, hence not in $M'$; let $\ast$ be the element of
$1(M)$. Then $\ast\in\hat G'(\neg s)$ and $\ast\in\cl_N(s)$. As $\neg
s\vee\neg\neg s$ is the top element, $\ast\in\hat G(\neg s)$, which
contradicts the homomorphism $a$, or $\ast\in\hat G(\neg\neg s)$,
which contradicts the homomorphism $\psi$. Hence $R'$ consists of the
$\iota a\psi^{-1}$.

Now we show $\hat G(U)=\hat G_a(U)$ for all models $\hat G$ with
non-empty $R$ and all $a\in R$ at once, by induction on the
construction of $U$; note that $\psi(\hat G_a(U))=\hat G_{a\psi^{-1}}
(U)$ for an isomorphism $\psi$ out of $M$, the two Kripke models being
isomorphic. The homomorphisms out of $\hat G$ are those with an
invertible underlying homomorphism, $\psi\colon\hat G\to\hat G'$ with
$\hat G'(U)=\hat G_{a\psi^{-1}}(U)$ by the induction hypothesis, among
them the identity, and the homomorphisms $\iota a$ into classical
models. So by Lemma~\ref{lem:truth}, $\alpha\in\hat G(\forall_fU)$ if
and only if $\beta\in\hat G_a(U)$ for all $\beta$ over $\alpha$ and
$\beta'\in\cl_N(U)$ for all $\beta'$ over $a(\alpha)$, which is the
forcing clause for $\forall_fU$ at the node $0$ of the Kripke model
of $a$. The case of $\Rightarrow$ is the same.
\end{proof}

\begin{thm}\label{thm:wem}
Let $L$ be an intermediate logic proving the weak excluded middle,
other than classical logic, and $\cC$ a pretopos with finite sorts
such that $\Mod(\cC)$ is an EI-category.
\begin{enumerate}
\item[(a)] If $\Mod(\cC)$ has no two composable non-invertible
homomorphisms, then $c^L_\cC$ is an equivalence.
\item[(b)] If $L$ does not prove linearity and $\Mod(\cC)$ has two
composable non-invertible homomorphisms, then $c^L_\cC$ is not an
equivalence.
\end{enumerate}
\end{thm}
\begin{proof}
(a) Let $L'$ be the logic of the Kripke frame $0<1$; it proves
linearity, $n_{L'}=2$, and $L\subseteq L'$, the frame $0<1$ being a
frame of $L$. By Lemma~\ref{lem:top} every model of $\HH_L\cC$ is the
evaluation, at a node, of the Heyting functor $\HH_L\cC\to\Set^K$ of a
Kripke model over the point or over the chain $0<1$; and $\Set^K$ is
an $L'$-category for these $K$. So every instance of a formula of
$L'$ has the top value in all models of $\HH_L\cC$, and is the top
subobject, the models jointly reflecting the order. Hence $\HH_L\cC$
is an $L'$-category, and it has the universal property of
$\HH_{L'}\cC$. By Theorem~\ref{thm:two}(b), $\HH_{L'}\cC$ is a
pretopos.

(b) Let $M\xrightarrow{h}N\xrightarrow{h'}N'$ be non-invertible. By
Proposition~\ref{prop:EI}(c) there are a sentence $s$ true in $N$ and
not in $M$, and a sentence $t$ true in $N'$ and not in $N$. The
diamond is a frame of $L$, and \cite[Prop.~5.1]{Interm} applies to the
terminal object $A=1$, the subobjects $S=s$ and $T=t$, the
homomorphisms $h$, $h'$ and the element $\ast$ of $1(M)$: we have
$\ast\notin s(M)$, $\ast\in s(N)\setminus t(N)$ and $\ast\in t(N')$.
So $c^L_\cC$ is not an equivalence.
\end{proof}

The Kripke model over the diamond used in \cite[Prop.~5.1]{Interm} has
$M$ at the root and at one of the middle nodes, $N$ at the other
middle node and $N'$ at the top:
\[
\begin{tikzcd}[column sep=small,row sep=small]
&N'&\\
N\ar[ur,"h'"]&&M\ar[ul,"h'h"']\\
&M\ar[ul,"h"]\ar[ur,"\id"']&
\end{tikzcd}
\]
The transition $\id$ is an isomorphism of the underlying models which
does not reflect the value of $s\Rightarrow t$.

\section{The classification and examples}

\begin{proof}[Proof of Theorem~\ref{thm:main}]
We may assume that $\cC$ is a pretopos (Section~\ref{sec:setting}).
(1) is Theorem~\ref{thm:two}(a). (2): $c^L_\cC$ is an equivalence if
and only if $\cC$ is Boolean \cite[Thm.~1.1(1)]{Interm}, that is, by
Lemma~\ref{lem:boolean}, if and only if $\cM$ is a groupoid. (4) is
Theorem~\ref{thm:two}(b). In (3) and (5) the logic is not classical,
and $\cM$ is an EI-category if $c^L_\cC$ is an equivalence, by the
remark after Theorem~\ref{thm:two}; the rest is
Theorem~\ref{thm:wem}.
\end{proof}

\begin{cor}[{\cite[Thm.~1.1]{Classify}}]
Let $\cX$ be a finite category in which idempotents split. For
$\cC=\Fin^{\cX}$ the statements of Theorem~\ref{thm:main} hold with
$\cM=\cX$.
\end{cor}

\begin{cor}[{\cite[Thm.~5.1]{Classify}}, {\cite[Thm.~6.3]{Interm}} for $n_L=2$]
Let $D$ be a bounded distributive lattice. If $L$ proves the weak
excluded middle and not linearity, $c^L_D$ is an equivalence if and
only if $D$ has no chain $P_0\subsetneq P_1\subsetneq P_2$ of prime
filters; if $L$ proves
linearity and $n_L=2$, it is an equivalence.
\end{cor}
\begin{proof}
The models of $D$ are the prime filters, a homomorphism is an
inclusion, and every object of $\PP D$ is a quotient of a finite
coproduct of subterminal objects; so $D$ has finite sorts and $\cM$
is a poset.
\end{proof}

\begin{ex}\label{ex:new}
A theory with finite sorts whose pretopos is neither of the form
$\Fin^{\cX}$ for a finite category $\cX$ nor that of a propositional
theory. Let $X$ be a non-empty finite poset, and $\bT$ the
coherent theory with a proposition $u_x$ for $x\in X$ and the axioms
\[
u_y\vdash u_x\ (x\le y),\qquad\top\vdash\bigvee_xu_x,\qquad u_x\wedge
u_y\vdash\bigvee_{z\ge x,y}u_z,
\]
saying that the true $u_x$ form the down-set of an element of $X$,
and with sorts $\sigma_e$, $e\in\mathbb N$, each with a relation $\ne$
and the axioms $x=y\wedge x\ne y\vdash\bot$, $\top\vdash x=y\vee x\ne
y$, $\top\vdash\exists x\exists y\,(x\ne y)$ and $\top\vdash x=y\vee
y=z\vee x=z$, saying that $\ne$ is the complement of the equality and
that $\sigma_e$ has exactly two elements. A model is an element of $X$
and a sequence of two-element sets, and there is a homomorphism from
the model of $x$ to that of $y$ iff $x\le y$; the category of models
is equivalent to $X\times\Gamma$ for the groupoid $\Gamma$ of the sequences of two-element sets
and bijections, an EI-category which is not equivalent to a finite
category, and not a poset. By Theorem~\ref{thm:main}, the translation
commutes with the imaginaries of $\bT$ for classical logic and for
the logic of the frame $0<1$; for the logics without the weak
excluded middle iff $X$ is discrete; and for those with the weak
excluded middle and without linearity iff $X$ has no chain of three
elements. For the logics proving linearity with $n_L\ge3$
Theorem~\ref{thm:main} decides only the case without chains of three
elements; but the theories in finitely many of the sorts $\sigma_e$
have the finite categories of models $X\times(\mathbb Z/2)^e$, up to
equivalence, the group being regarded as a category with one object, for which the comparison is an equivalence
\cite[Cor.~6.2]{Strings}, and the pretopos of $\bT$ is the directed
colimit of theirs; so by Lemma~\ref{lem:colimit} the translation
commutes with the imaginaries of $\bT$ for every logic proving
linearity.
\end{ex}

\begin{ex}
Theories with a coherent inequality. Let $\bT$ be a coherent theory in
which every sort has a relation $\ne$ which is the complement of the
equality, and suppose that all models are finite on every sort. Every
homomorphism is injective on the sorts of the language, and an injective endomorphism of a model
with finite sorts is bijective on every sort and reflects the
relations, by counting; so $\cM$ is an EI-category, and case (4) of
Theorem~\ref{thm:main} applies: the translation into the logic of the
frame $0<1$ commutes with the imaginaries. An example is the theory of
a sort with at most $n$ elements and a relation $\ne$, whose models
are the sets with at most $n$ elements and the injective maps.
\end{ex}

\section{Questions}

\begin{que}
For the logics proving linearity with $n_L\ge3$ and a pretopos $\cC$
with finite sorts whose category of models is an EI-category: is
$c^L_\cC$ an equivalence if and only if every finite language is
contained in a finite language $T$ such that $\cX_T$ is an
EI-category without obstruction in the sense of
\cite[Thm.~6.15]{Strings}? The condition is sufficient by
Lemmas~\ref{lem:union} and~\ref{lem:colimit}.
\end{que}

\begin{que}
Is the absence of homomorphisms $i$ and objects $Y$ as in
Theorem~\ref{thm:obstruction}, together with the absence of
configurations as in \cite[Prop.~5.1]{Interm} where the diamond is a
frame, sufficient for $c^L_\cC$ to be an equivalence for some class of
theories with infinite models? For the theory of pure equality the comparison is
not an equivalence for any non-classical logic \cite{Interm}.
\end{que}

\section*{Declarations}
The author has no competing interests to declare.

\begin{thebibliography}{9}
\bibitem{Heyting} R.~Shibusawa, \emph{The free Heyting category on a
pretopos is a pretopos only in the Boolean case: the intuitionistic
translation and the pretopos completion}, preprint (2026), archived at
DOI 10.5281/zenodo.24182122.
\bibitem{Interm} R.~Shibusawa, \emph{Intermediate logics and
imaginaries: the translation into an intermediate logic commutes with
the pretopos completion of all finite propositional theories exactly
for the extensions of G\"odel--Dummett logic}, preprint (2026),
archived at DOI 10.5281/zenodo.24150610.
\bibitem{Morley} R.~Shibusawa, \emph{Translations in the lattice of
two-out-of-three classes: Morleyization as a sameness of theories, its
join with Morita equivalence, and the forced order of translation and
completion}, preprint (2026), archived at DOI 10.5281/zenodo.24122457.
\bibitem{Strings} R.~Shibusawa, \emph{G\"odel--Dummett logic over a
finite category of models: the models of the translated theory are
strings of arrows, and a subgroup obstruction to commuting with the
imaginaries}, preprint (2026), archived at DOI
10.5281/zenodo.24327653.
\bibitem{Classify} R.~Shibusawa, \emph{Intermediate logics and the
imaginaries of a theory with a finite category of models: a
classification}, preprint (2026), archived at DOI
10.5281/zenodo.24331748.
\bibitem{Ye24} L.~Ye, \emph{Stack representation of finitely presented
Heyting pretoposes I}, preprint, arXiv:2402.13099.
\bibitem{Ye26} L.~Ye, \emph{Duality theory for categorical theories},
preprint, arXiv:2605.22071.
\end{thebibliography}

\end{document}
